% ------------------------------------------------------------------------------
% Chapter 2
% ------------------------------------------------------------------------------
\chapter{Literature Review} % enter the name of the chapter here
\label{cha:ch2_literature} % enter the chapter label here (for cross referencing)

% ------------------------------------------------------------------------------
% Introduction
% ------------------------------------------------------------------------------
\section{Introduction}
\label{sec:Introduction}
Generative AI has become one of the most used technologies in different fields to perform various tasks where high levels of knowledge and information are required. This is different from traditional AI in the sense that this technology emphasizes cognitive complexity within the output created and responding to the need for community-specific interpretation, refinement, and integration. In software development, artificial intelligence is used in the form of coding-monitoring assistant tools that are included in the software-development environment to facilitate code generation, fixing bugs in the software, and creating documentation.

Empirical studies on the impact of generative AI on human productivity and output quality revealed that, under conditions of structured task design, a significant increase in human output is achieved. In a large-scale experimental design, the availability of generative AI tools resulted in a reduction on task completion time and an increase in output quality. The results achieved were statistically significant for the total sample, while the results were more pronounced for low-performing individuals \autocite{Noy2023}. The results support the view that generative AI tools can produce productivity gains at the individual task level.

However, the productivity gains at the task level do not necessarily translate into organizational return on investment. Research on AI investment at the firm level indicates that the performance effects are heterogeneous and conditional. AI investment is linked with growth through innovation and product development channels; however, the level and stability of the returns depend on the capabilities of the organization \autocite{Babina2024}. The market reaction to AI investment announcements is mixed, indicating the level of awareness among investors about the risk and capabilities \autocite{Lui2021}.

These studies all circle back to the same basic idea: AI capability does not produce value; rather, value is produced when the technological tools are embedded in the organizational systems. More recent studies have suggested another component: a digital organizational culture can influence the innovation outcome of AI capability \autocite{An2024}. Leadership can also play a part in influencing the development of trust and embedding AI usage \autocite{Schneider2023}.

Despite the recent increase in research on the performance effects of AI, little research has appeared that combines the effects of organizational culture and leadership into a structured ROI model for AI coding assistants in software development teams. This chapter synthesizes global empirical research on AI productivity and firm performance, critically evaluates the impact of organizational conditions, and focuses the research question more specifically on the South African market, where empirical ROI research has yet to be conducted.

